\chapter{THẾ HỆ MỚI: YOLOv9 -- YOLOv12 (2024--2025)}

%% ========== YOLOv9 ==========
\section{YOLOv9 (2024) --- Chống mất thông tin}
\textbf{Paper:} ``YOLOv9: Learning What You Want to Learn Using Programmable Gradient Information''\\
\textbf{Tác giả:} Chien-Yao Wang, I-Hau Yeh, Hong-Yuan Mark Liao

\subsection{Vấn đề: Information Bottleneck}
Khi data đi qua nhiều layers, \textbf{mutual information giảm đơn điệu}:
\[I(X,X) \geq I(X, f_1(X)) \geq I(X, f_2(f_1(X))) \geq \ldots\]

\textit{Ví dụ:} Giống trò ``truyền tin'' --- người đầu nhận tin gốc, truyền miệng qua 20 người, người cuối nhận tin sai lệch. Layers sâu ``quên'' thông tin input gốc $\rightarrow$ gradient không đáng tin cậy $\rightarrow$ model không học tối ưu.

\subsection{PGI (Programmable Gradient Information)}
3 thành phần:
\begin{enumerate}
    \item \textbf{Main Branch:} Network chính dùng cho inference. Kiến trúc tùy chọn (không bị ràng buộc).
    \item \textbf{Auxiliary Reversible Branch:} Mạng phụ có khả năng \textbf{khôi phục input từ output} $\rightarrow$ giữ complete information $\rightarrow$ gradient đáng tin cậy. \textbf{Chỉ dùng khi training, bỏ khi inference} $\rightarrow$ zero inference cost.
    \item \textbf{Multi-level Auxiliary:} Inject gradient ở \textbf{nhiều tầng} --- tránh vấn đề gradient bị ``đứt'' khi main branch rất sâu.
\end{enumerate}

\begin{analogybox}
\textbf{Truyền tin và SMS:} Main branch giống ``truyền miệng'' qua 20 người (mất thông tin). PGI thêm 1 \textbf{kênh SMS song song} (auxiliary branch) --- mỗi người đều nhận SMS đối chiếu với tin truyền miệng $\rightarrow$ sửa sai kịp thời.
\end{analogybox}

\textbf{Tại sao PGI tốt hơn Deep Supervision?}
\begin{itemize}
    \item Deep supervision: Thêm loss ở giữa, nhưng features ở giữa \textbf{đã mất thông tin} $\rightarrow$ gradient dựa trên info không đầy đủ
    \item PGI: Auxiliary reversible branch \textbf{giữ full information} $\rightarrow$ gradient luôn đáng tin cậy
    \item PGI đặc biệt hiệu quả cho \textbf{lightweight models}: cải thiện \textbf{+2.5\% AP} (vs +0.9\% cho medium model)
\end{itemize}

\subsection{GELAN (Generalized ELAN)}
Tổng quát hóa CSPNet + ELAN, cho phép dùng \textbf{bất kỳ computational block} nào bên trong (Conv, Bottleneck, Transformer...).

\textit{Ví dụ:} ELAN cũ giống dây chuyền chỉ dùng được 1 loại máy. GELAN giống dây chuyền ``universal'' --- lắp máy gì vào cũng chạy được.

Thí nghiệm thay block:
\begin{table}[H]\centering
\caption{GELAN --- tính linh hoạt khi thay block}
\begin{tabular}{lc}\toprule
\textbf{Block} & \textbf{mAP50-95}\\\midrule
ELAN + CSPBlock & 52.3\% \\
ELAN + DarkBlock & 52.4\% \\
\textbf{GELAN (mixed)} & \textbf{53.0\%} \\\bottomrule
\end{tabular}
\end{table}

\subsection{Kết quả chi tiết}
\begin{table}[H]\centering
\caption{YOLOv9 so sánh SOTA (COCO val2017)}
\begin{tabular}{lcccl}\toprule
\textbf{Model} & \textbf{Params} & \textbf{FLOPs} & \textbf{mAP50-95} & \textbf{So với v8}\\\midrule
YOLOv8-S & 11.2M & 28.6G & 44.9\% & --- \\
\textbf{YOLOv9-S} & \textbf{7.2M} & \textbf{26.7G} & \textbf{46.8\%} & +1.9\%, ít params hơn \\
YOLOv8-M & 25.9M & 78.9G & 50.2\% & --- \\
\textbf{YOLOv9-C} & \textbf{25.5M} & \textbf{102.8G} & \textbf{53.0\%} & +2.8\% \\
YOLOv8-X & 68.2M & 257.8G & 53.9\% & --- \\
\textbf{YOLOv9-E} & \textbf{58.1M} & \textbf{192.5G} & \textbf{55.6\%} & +1.7\%, ít params hơn \\\bottomrule
\end{tabular}
\end{table}

\textbf{Nhận xét:} v9 luôn \textbf{ít params hơn} v8 mà accuracy cao hơn ở mọi scale --- chứng minh PGI hiệu quả.

%% ========== YOLOv10 ==========
\section{YOLOv10 (2024) --- NMS-Free}
\textbf{Paper:} ``YOLOv10: Real-Time End-to-End Object Detection''\\
\textbf{Tác giả:} Ao Wang et al. (Tsinghua University)

\subsection{Vấn đề NMS}
NMS (Non-Maximum Suppression) có 3 nhược điểm:
\begin{enumerate}
    \item \textbf{Tốn thời gian:} Cần tune confidence threshold, IoU threshold
    \item \textbf{Không differentiable:} Không tối ưu end-to-end được
    \item \textbf{Latency phụ thuộc data:} Ảnh nhiều vật $\rightarrow$ NMS chậm hơn $\rightarrow$ latency không ổn định
\end{enumerate}

\begin{analogybox}
\textbf{Dây chuyền zero defect:} NMS giống kiểm tra chất lượng cuối dây chuyền --- phải dừng lại soi từng sản phẩm, loại bỏ phế phẩm. v10 thiết kế dây chuyền ``zero defect'' --- không cần kiểm tra cuối.
\end{analogybox}

\subsection{Consistent Dual Assignments}
\begin{itemize}
    \item \textbf{One-to-many (o2m):} 1 vật $\rightarrow$ nhiều predictions $\rightarrow$ nhiều training signal $\rightarrow$ train tốt
    \item \textbf{One-to-one (o2o):} 1 vật $\rightarrow$ 1 prediction $\rightarrow$ không cần NMS $\rightarrow$ deploy tốt
    \item \textbf{Dual heads --- giải pháp:}
    \begin{itemize}
        \item Training: Dùng \textbf{cả 2 heads} --- o2m tạo dense supervision, o2o học matching
        \item Inference: Chỉ dùng \textbf{o2o head} $\rightarrow$ không cần NMS
    \end{itemize}
\end{itemize}

\subsection{Efficiency-Driven Design}
\begin{table}[H]\centering
\caption{5 kỹ thuật efficiency của YOLOv10}
\begin{tabular}{lll}\toprule
\textbf{Kỹ thuật} & \textbf{Ý tưởng} & \textbf{Hiệu quả}\\\midrule
Lightweight cls head & Depthwise separable conv & Giảm computation \\
Decoupled downsample & Tách spatial + channel & Rẻ hơn \\
Rank-guided block & Block nhẹ ở stage low-rank & Tự động chọn \\
Large-kernel conv & $7\times7$ depthwise ở stages sâu & Tăng receptive field \\
PSA & 50\% channels qua attention & Giảm cost 50\% \\\bottomrule
\end{tabular}
\end{table}

\subsection{Kết quả}
\begin{table}[H]\centering
\caption{YOLOv10 --- NMS-free so sánh}
\begin{tabular}{lccccl}\toprule
\textbf{Model} & \textbf{Params} & \textbf{Latency} & \textbf{mAP} & \textbf{So sánh}\\\midrule
YOLOv8-S + NMS & 11.2M & 4.70ms & 44.9\% & Cần NMS \\
o2o only training & --- & 2.33ms & 43.8\% & Kém accuracy \\
\textbf{v10-S Dual} & \textbf{7.2M} & \textbf{2.49ms} & \textbf{46.3\%} & NMS-free, nhanh 2$\times$ \\\midrule
YOLOv8-X & 68.2M & 7.61ms & 53.9\% & --- \\
\textbf{v10-X} & \textbf{29.5M} & \textbf{10.70ms} & \textbf{54.4\%} & 2.3$\times$ ít params \\\bottomrule
\end{tabular}
\end{table}

\textbf{Trade-off:} v10-X (54.4\%) \textbf{thua} v9-E (55.6\%) về mAP. Nhưng v10 ít params gần \textbf{2$\times$} và \textbf{không cần NMS}. v10 đổi accuracy lấy efficiency.

%% ========== YOLO11 ==========
\section{YOLO11 (2024) --- Nâng cấp toàn diện}
\textbf{Tác giả:} Ultralytics \hfill \textit{Không có paper} \hfill \textit{Lưu ý: bỏ ``v'' --- YOLO11 thay vì YOLOv11}

\subsection{C3K2 (CSP Bottleneck with 2 Kernel sizes)}
Kế thừa C2f (v8) nhưng dùng \textbf{2 kernel sizes: $3\times3$ và $5\times5$}:
\begin{itemize}
    \item Kernel $3\times3$: Capture \textbf{local features} (chi tiết nhỏ, cạnh, góc)
    \item Kernel $5\times5$: Capture \textbf{wider context} (vùng rộng hơn)
    \item Multi-scale feature extraction \textbf{trong cùng 1 block} $\rightarrow$ phong phú hơn C2f (chỉ dùng $3\times3$)
\end{itemize}

\textit{Ví dụ:} C2f giống nhìn qua 1 ống kính cố định. C3K2 giống nhìn qua 2 ống kính zoom khác nhau cùng lúc --- vừa thấy chi tiết, vừa thấy bối cảnh.

\subsection{C2PSA (Cross Stage Partial Spatial Attention)}
Thêm \textbf{spatial attention} vào CSP:
\begin{enumerate}
    \item Features đi qua CSP split (2 nhánh)
    \item Một nhánh qua \textbf{spatial attention module} --- học ``vùng nào quan trọng''
    \item Attention map highlight vùng có vật thể, suppress background
    \item Concat với nhánh bypass $\rightarrow$ output
\end{enumerate}

\textit{Ví dụ:} C2f giống quét đều toàn bộ ảnh (mỗi pixel được xử lý ngang nhau). C2PSA giống mắt người --- \textbf{tập trung vào vùng quan trọng}, lướt nhanh qua background.

\subsection{Kết quả}
\begin{table}[H]\centering
\caption{YOLO11 so sánh v8 (COCO val2017, input 640)}
\begin{tabular}{lcccl}\toprule
\textbf{Model} & \textbf{Params} & \textbf{FLOPs} & \textbf{mAP50-95} & \textbf{vs v8}\\\midrule
11n & 2.6M & 6.5G & 39.5\% & v8n: 37.3\% (\textbf{+2.2\%}) \\
11s & 9.4M & 21.5G & 47.0\% & v8s: 44.9\% (\textbf{+2.1\%}) \\
11m & 20.1M & 68.0G & 51.5\% & v8m: 50.2\% (\textbf{+1.3\%}) \\
11l & 25.3M & 87.6G & 53.4\% & v8l: 52.9\% (\textbf{+0.5\%}) \\
11x & 56.9M & 194.9G & 54.7\% & v8x: 53.9\% (\textbf{+0.8\%}) \\\bottomrule
\end{tabular}
\end{table}

\textbf{Nhận xét:} Cải thiện ở \textbf{mọi scale}, đặc biệt mạnh ở Nano/Small (+2\%). Params tương đương hoặc ít hơn v8.

%% ========== YOLOv12 ==========
\section{YOLOv12 (2025) --- Kỷ nguyên Attention}
\textbf{Paper:} ``YOLOv12: Attention-Centric Real-Time Object Detectors''\\
\textbf{Tác giả:} Yunjie Tian, Qixiang Ye, David Doermann

\subsection{Thách thức đưa Attention vào YOLO}
\begin{enumerate}
    \item \textbf{Tốc độ:} Self-attention có complexity $O((H\times W)^2)$ --- quá chậm cho real-time
    \item \textbf{Optimization:} Attention networks khó train hơn CNN (gradient không ổn định)
    \item \textbf{Hardware:} Attention operations chưa tối ưu trên GPU bằng Conv
\end{enumerate}

\textit{Ví dụ:} CNN giống đọc sách từng dòng (local, nhanh). Attention giống nhìn cả trang cùng lúc (global, hiểu sâu nhưng chậm). v12 tìm cách ``nhìn cả trang'' mà vẫn nhanh.

\subsection{Area Attention (A$^2$)}
Chia feature map thành \textbf{$n$ vùng bằng nhau}, tính attention \textbf{trong từng vùng riêng biệt}:
\[O(\text{Global}) = (H \times W)^2 \quad \rightarrow \quad O(\text{Area, } n=4) = H \times W \times \frac{HW}{4}\]

Giảm \textbf{4$\times$ complexity}, chỉ mất \textbf{0.1\% AP}.

3 cách chia vùng: ngang (n vùng ngang), dọc (n vùng dọc), grid ($\sqrt{n} \times \sqrt{n}$).

\begin{analogybox}
\textbf{Lớp học chia nhóm:} Global attention giống 1 giáo viên kiểm tra bài của \textbf{tất cả} 100 học sinh cùng lúc (quá tải). Area attention chia thành 4 nhóm 25 người, kiểm tra từng nhóm $\rightarrow$ nhanh 4$\times$, kết quả gần tương đương.
\end{analogybox}

\subsection{R-ELAN (Residual ELAN)}
\textbf{Vấn đề:} Thay CNN bằng attention trong ELAN $\rightarrow$ training bất ổn, accuracy giảm. Nguyên nhân: attention outputs có \textbf{variance lớn} $\rightarrow$ tích lũy khi concat $\rightarrow$ gradient explode.

\textbf{Giải pháp:} Thêm residual shortcut với \textbf{scale factor 0.01}:
\begin{itemize}
    \item Scale nhỏ $\rightarrow$ output gần identity ở đầu training $\rightarrow$ gradient ổn định
    \item Dần dần model học scale factor phù hợp
    \item ELAN + Attention: 39.1\% (bất ổn) $\rightarrow$ \textbf{R-ELAN + Attention: 40.6\% (ổn định)}
\end{itemize}

\subsection{Tối ưu phần cứng}
\begin{itemize}
    \item \textbf{FlashAttention:} Giảm memory footprint cho attention
    \item \textbf{Conv+BN thay Linear+LN:} Tận dụng GPU parallelism tốt hơn
    \item \textbf{MLP ratio giảm:} 4.0 (Transformer chuẩn) $\rightarrow$ 1.2--2.0
    \item \textbf{Bỏ Positional Encoding:} Thay bằng $7\times7$ depthwise conv
\end{itemize}

\subsection{Kết quả chi tiết}
\begin{table}[H]\centering
\caption{YOLOv12 so sánh (COCO val2017)}
\begin{tabular}{lcccc}\toprule
\textbf{Model} & \textbf{Params} & \textbf{GFLOPs} & \textbf{mAP50-95}\\\midrule
YOLOv10-N & 2.3M & 6.7 & 38.5\% \\
YOLO11-N & 2.6M & 6.5 & 39.5\% \\
\textbf{YOLOv12-N} & \textbf{2.6M} & \textbf{6.5} & \textbf{40.6\%} \\
\midrule
YOLO11-S & 9.4M & 21.5 & 47.0\% \\
\textbf{YOLOv12-S} & \textbf{9.3M} & \textbf{21.4} & \textbf{48.0\%} \\
\midrule
YOLO11-L & 25.3M & 87.6 & 53.4\% \\
\textbf{YOLOv12-L} & \textbf{26.4M} & \textbf{88.9} & \textbf{53.7\%} \\
\midrule
\textbf{YOLOv12-X} & \textbf{59.1M} & \textbf{199.0} & \textbf{55.2\%} \\\bottomrule
\end{tabular}
\end{table}

\textbf{Ý nghĩa:} YOLOv12 đánh dấu YOLO chuyển từ \textbf{CNN-centric} sang \textbf{Attention-centric}. Nhược điểm: cần GPU hỗ trợ FlashAttention (GPU hiện đại).
